\appendix

\section{Training Configurations}
\label{app:training-config}

Table~\ref{tab:primary-config} gives the relational settings saved
with the four seed-0 checkpoints. Onset is the fraction of completed
optimization steps. Standard MLM prediction always uses the complete
384-dimensional state. The two single-target systems also apply their
relational loss to the complete state; only Relay Joint uses fixed
coordinate slices.

\begin{table}[H]
\centering
\small
\setlength{\tabcolsep}{2pt}
\begin{tabular}{@{}p{0.13\textwidth}p{0.15\textwidth}p{0.17\textwidth}p{0.17\textwidth}p{0.18\textwidth}p{0.12\textwidth}@{}}
\toprule
Method & Syn onset / weight & Para syntax onset / weight & Para content onset / weight & Relational rep. & Entity anchor \\
\midrule
Standard MLM & Off & Off & Off & -- & Off \\
Syntagmatic & $0.0920386562$ / $5\!\times\!10^{-4}$ & Off & Off & Full 384 & Off \\
Paradigmatic & Off & $0.5522310170$ / $5\!\times\!10^{-5}$ & $0.99$ / $10^{-6}$ & Full 384 & Off \\
Relay Joint & $0.0920386562$ / $5\!\times\!10^{-4}$ & $0.1840773125$ / $5\!\times\!10^{-5}$ & $0.99$ / $10^{-6}$ & Syn 0--255; Para 256--383 & $\approx9.2\%$--$60\%$ / $2.5\!\times\!10^{-4}$ \\
\bottomrule
\end{tabular}
\caption{Saved relational configurations. After Relay Joint activates
syntax Paradigmatic supervision, its Syntagmatic weight at eligible
syntax positions is retained at 10\% of the earlier value.}
\label{tab:primary-config}
\end{table}

Both targets use $K=8$ and temperature 1.0. Syntagmatic takes
attention from the last layer. Paradigmatic has zero weight for entity
and reading-sensitive categories and no post-onset ramp. Relay Joint
does not enable a content projection or learned router. Its routing is
a fixed coordinate selection, and its anchor has no trainable
parameters. The total parameter count is 34,677,952 for every system.

\begin{table}[H]
\centering
\small
\setlength{\tabcolsep}{4pt}
\begin{tabular}{p{0.20\textwidth}p{0.72\textwidth}}
\toprule
Objective & Eligible MLM-selected positions \\
\midrule
Syntagmatic & Gold-token category is syntax, reading, or content. Entity and invalid tokens are excluded. \\
Paradigmatic, syntax & Gold token matches the fixed syntax list. It begins at progress $0.5522310170$ in Paradigmatic and $0.1840773125$ in Relay Joint, with weight $5\times10^{-5}$. \\
Paradigmatic, content & Gold token is classified as content; active after progress $0.99$ with weight $10^{-6}$. Entity and reading categories receive zero Paradigmatic weight. \\
Entity anchor & Relay Joint only: gold token is classified as entity; active from progress $0.0920386562$ through $0.60$ with weight $2.5\times10^{-4}$. \\
\bottomrule
\end{tabular}
\caption{Eligibility for the relational objectives and Relay Joint
entity anchor. Categories are fixed gold-token ID lookup tables
constructed from the tokenizer vocabulary before training.}
\label{tab:token-eligibility}
\end{table}

\paragraph{Vocabulary-derived token categories.}
No POS tagger or NER system is used. The tokenizer vocabulary is
scanned once before training. One recognized leading SentencePiece,
byte-level BPE, WordPiece, or space marker is removed; trailing
punctuation is stripped; and fixed-list matching is case-insensitive.
Byte tokens and entries with no alphabetic character are invalid.
The lookup then assigns the first matching category in this order:
syntax, reading-sensitive, entity, and content. Thus the categories are
mutually exclusive even when lexical lists overlap. An entity is a
remaining entry containing a digit, or an alphabetic entry of at least
three characters that begins with an uppercase letter. Remaining
alphabetic entries are content tokens.

\paragraph{Fixed lexical lists.}
\textbf{Auxiliaries:} is, are, was, were, be, been, being, am, do, does,
did, have, has, had, can, could, will, would, should, may, might, must,
shall, ought, need, dare, used.
\textbf{Determiners:} a, an, the, this, that, these, those, some, any,
all, each, every, many, few, no, much, more, most, several, both,
either, neither, such, what.
\textbf{Pronouns:} he, she, it, they, him, her, them, his, their, its,
we, us, our, you, your, i, me, my, myself, yourself, himself, herself,
itself, ourselves, themselves, one, ones.
\textbf{Negation:} not, n\textquotesingle t, never, no, nor, neither.
\textbf{Prepositions:} in, on, at, by, with, from, to, of, for, into,
onto, over, under, near, about, between, through, during, without,
within, along, across, behind, beyond, toward, towards, upon, among,
amongst, beside, besides, against, around, before, after, above, below,
off, up, down, out.
\textbf{Conjunctions:} and, or, but, because, although, if, when, while,
before, after, since, until, unless, whereas, so, yet, than, as, though,
whether, once, till, lest, except, provided, given, suppose, assuming,
whenever, wherever.
\textbf{Wh-words:} who, whom, whose, which, where, why, how, what,
whatever, whichever, whoever, however, wherever.
\textbf{Comparatives:} more, less, fewer, better, worse, bigger,
smaller, higher, lower, longer, shorter, older, younger.
\textbf{Temporal expressions:} now, then, ago, later, earlier, soon,
already, still, yet, finally, eventually, previously, formerly,
currently, recently, lately, immediately, suddenly, gradually.

\paragraph{Gold-token and context handling.}
Target-position eligibility is looked up from the clean gold/original
token ID, never the corrupted input. The same category therefore
applies whether an MLM-selected token is replaced with \texttt{[MASK]},
a random token, or left unchanged. The Syntagmatic context set excludes
padding, \texttt{[CLS]}, \texttt{[SEP]}, \texttt{[MASK]}, the target
itself, and every other MLM-selected position. Context embeddings are
then gathered from the original uncorrupted token IDs. Subword entries
are classified independently; the code performs no word reconstruction.

\section{Reproduction Mapping}
\label{app:reproduction}

Table~\ref{tab:model-mapping} maps each paper name to the exact seed-0
model directory used in the main comparison. All reported final scores
use the \texttt{chck\_100M} checkpoint inside that directory.

\begin{table}[H]
\centering
\small
\begin{tabular}{lp{0.70\textwidth}}
\toprule
Paper name & Exact model directory \\
\midrule
Standard MLM & \texttt{bb26\_claim\_clean\_s0} \\
Syntagmatic & \texttt{bb26\_claim\_comp\_entguard\_s0} \\
Paradigmatic & \texttt{bb26\_claim\_agg\_synonly\_s0} \\
Relay Joint & \texttt{bb26\_claim\_compagg\_v22\_v15\_relay\_synp050\_synkeep010\_s0} \\
\bottomrule
\end{tabular}
\caption{Paper-to-directory mapping for the seed-0 comparison.}
\label{tab:model-mapping}
\end{table}

Each model directory contains its saved training command, trainer
copy, and claim manifest. Evaluation results use the BabyLM evaluation
pipeline at commit
\texttt{3d57ddc8c40ee795c0b5e41b3a20251a9457a593}. The local evaluation
repository contained uncommitted changes when this version was
audited, so the hash identifies its HEAD rather than a clean checkout.

Figure~\ref{fig:reading-aoa} is generated from the official checkpoint
results for the four directories above. Figure
\ref{fig:contextual-alignment} uses the exact 100M checkpoints and the
new seed-0 CSVs listed in the accompanying v22 analysis directory.

\section{Representation Analyses}
\label{app:representation-analysis}

\paragraph{Paradigmatic neighborhood probe.}
The probe uses 2,000 eligible masked syntax-token instances. For each
instance, the correct candidate is the centroid of the eight nearest
input-embedding neighbors of the gold token. Sixty-three centroids from
other syntax-token IDs form controls. Alignment margin measures the
cosine advantage of the correct target over controls; MRR is its
reciprocal rank. Table~\ref{tab:para-probe} reports paired item-level
differences from the exact seed-0 Standard MLM checkpoint. Full-384 is
the primary comparison; equal-width slices are diagnostic coordinate
views.

\begin{table}[H]
\centering
\small
\begin{tabular}{lrr}
\toprule
Method / representation & Margin $\Delta$ & MRR $\Delta$ \\
\midrule
Paradigmatic / Full-384 & .024213 & .084507 \\
Relay Joint / Full-384 & .047681 & .144508 \\
Relay Joint / Slice 1 (0--127) & .125988 & .218807 \\
Relay Joint / Slice 2 (128--255) & -.082908 & -.046663 \\
Relay Joint / Slice 3 (256--383) & .087723 & .112787 \\
\bottomrule
\end{tabular}
\caption{Seed-0 neighborhood-alignment differences from Standard MLM.
No cross-seed standard deviation or significance test is available.}
\label{tab:para-probe}
\end{table}

The final 128 coordinates are the fixed Paradigmatic training slice
for Relay Joint. Their positive diagnostic values show alignment in
that designated coordinate range, but do not imply that the model
discovered an intrinsic linguistic subspace.

\paragraph{Syntagmatic diagnostics.}
Attention is ranked over SentencePiece positions before word-level
mapping. Selected subwords are mapped to Universal Dependencies words,
with duplicate words collapsed. Recall@8 is the fraction of gold direct
dependency-neighbor words recovered; Precision@8 is the fraction of
unique mapped selected words that are direct dependency neighbors. For
context removal, the same selected token positions are zeroed in the
input to the final Transformer layer. A matched control removes the
same number of random non-special, non-target positions. Tables
\ref{tab:dependency} and~\ref{tab:context-removal} include only the
current Standard MLM and Syntagmatic runs.

\begin{table}[H]
\centering
\small
\begin{tabular}{llrr}
\toprule
Method & Seed / summary & Recall@8 & Precision@8 \\
\midrule
Standard MLM & 0 / 1 / 42 & .8707 / .8806 / .8699 & .2997 / .3003 / .2970 \\
 & Mean $\pm$ SD & $.8737\pm.0059$ & $.2990\pm.0018$ \\
Syntagmatic & 0 / 1 / 42 & .8750 / .8759 / .8772 & .2999 / .2979 / .2987 \\
 & Mean $\pm$ SD & $.8760\pm.0011$ & $.2988\pm.0010$ \\
\bottomrule
\end{tabular}
\caption{Dependency alignment by seed and sample mean $\pm$ SD.}
\label{tab:dependency}
\end{table}

\begin{table}[H]
\centering
\small
\begin{tabular}{llrrr}
\toprule
Method & Seed / summary & Attended & Matched & Difference \\
\midrule
Standard MLM & 0 / 1 / 42 & 1.6525 / 1.6137 / .9038 & -.0447 / .1678 / -.0039 & 1.6972 / 1.4459 / .9077 \\
 & Mean $\pm$ SD & $1.3900\pm.4215$ & $.0397\pm.1128$ & $1.3503\pm.4034$ \\
Syntagmatic & 0 / 1 / 42 & 1.8058 / 2.1002 / .7960 & -.0438 / .0159 / -.0082 & 1.8496 / 2.0842 / .8042 \\
 & Mean $\pm$ SD & $1.5673\pm.6840$ & $-.0120\pm.0300$ & $1.5793\pm.6814$ \\
\bottomrule
\end{tabular}
\caption{Gold-token NLL increase after context removal, by seed and sample mean $\pm$ SD. Difference is attended minus matched.}
\label{tab:context-removal}
\end{table}
